\chapter{Теория}

\section{Логическое И}

\begin{figure}[H] 
    \centering
    \includegraphics[width=0.75\textwidth]{table.png}
    \caption{Таблица истинности для логического И}
\end{figure}

\section{Сравнение IP-адреса источника и назначения}

Имеется сеть с такой топологией

\begin{figure}[H] 
    \centering
    \includegraphics[width=0.75\textwidth]{nw1.png}
    \caption{Топология сети}
\end{figure}

Перед отправкой пакета PC-1 узнает, в какой сети он находится. Для
этого PC-1 побитно произведет логическую операцию «И» между своим IP-адресом и своей
маской подсети

\begin{figure}[H] 
    \centering
    \includegraphics[width=0.75\textwidth]{and1.png}
    \caption{Определение своей сети}
\end{figure}

После того, как PC-1 определил, в какой сети находится, он побитно совершает логическое
«И» над IP-адресом назначения со своей маской и сравнивает с ранее полученным результатом

\begin{figure}[H] 
    \centering
    \includegraphics[width=0.75\textwidth]{and2.png}
    \caption{Определение сети получателя}
\end{figure}

Как видно, результат одинаковый, а это значит, что PC-2 со своим IP-адресом находится в
одной сети с PC-1.




\endinput